\documentclass{article}
\usepackage[utf8]{inputenc}
\usepackage{hyperref}
\usepackage{listings}
\usepackage{xcolor}
\usepackage{enumitem}

\lstset{
    basicstyle=\ttfamily\small,
    breaklines=true,
    frame=single,
    backgroundcolor=\color{gray!5},
}

\title{Load CSV file and print dataset statistics using R}
\author{marcos}
\date{2026-04-25}

\begin{document}

\section*{Load CSV file and print dataset statistics using R}
\textbf{Author:} marcos \hfill \textbf{Date:} 2026-04-25

\noindent Loads a CSV file using R's read.csv and prints rows, columns, total cells, file size, and object memory usage.

\section*{Language}
rscript

\section*{Category}
data-processing

\section*{Command}
\begin{lstlisting}
Rscript -e "args<-commandArgs(trailingOnly=TRUE);if(length(args)<1)stop('Need filename');df<-read.csv(args[1]);cat('Rows:',nrow(df),'Cols:',ncol(df),'Cells:',nrow(df)*ncol(df),'File size:',file.size(args[1]),'Object size:',format(object.size(df),units='auto'),'\n')"
\end{lstlisting}

\section*{Explanation}
The command uses Rscript to execute R code that reads a CSV file specified as argument, computes basic statistics (rows, columns, cells), gets file size via file.size(), and reports object memory usage via object.size().

\section*{Tags}
rscript, csv, statistics, data-analysis, memory

\section*{Dependencies}
r-base-core


\section*{Arguments}
\begin{enumerate}
  \item \textbf{CSV\_FILE} (Required): CSV file to analyze
\end{enumerate}

\section*{Examples}
\begin{enumerate}
  \item \texttt{Rscript -e "args<-commandArgs(trailingOnly=TRUE);if(length(args)<1)stop('Need filename');df<-read.csv(args[1]);cat('Rows:',nrow(df),'Cols:',ncol(df),'Cells:',nrow(df)*ncol(df),'File size:',file.size(args[1]),'Object size:',format(object.size(df),units='auto'),'\textbackslash\{\}n')" data.csv} - Analyze data.csv
  \item \texttt{Rscript -e "args<-commandArgs(trailingOnly=TRUE);if(length(args)<1)stop('Need filename');df<-read.csv(args[1]);cat('Rows:',nrow(df),'Cols:',ncol(df),'Cells:',nrow(df)*ncol(df),'File size:',file.size(args[1]),'Object size:',format(object.size(df),units='auto'),'\textbackslash\{\}n')" large\_dataset.csv 2>/dev/null} - Analyze suppressing warnings
\end{enumerate}

\section*{Output}
\begin{lstlisting}
Rows: 100 Cols: 10 Cells: 1000 File size: 12345 Object size: 78.2 Kb

\end{lstlisting}

\section*{Notes}
\begin{itemize}
  \item Uses read.csv with default settings (header=TRUE, stringsAsFactors=FALSE in R>=4.0)
  \item File size reported in bytes
  \item Object size includes R overhead
  \item For large files, consider adding options(stringsAsFactors=FALSE)
\end{itemize}

\section*{Warnings}
\begin{itemize}
  \item May load entire file into memory
  \item Very large CSV files could exhaust memory
\end{itemize}

\section*{See Also}
\begin{itemize}
  \item tsv-stats
\end{itemize}

\section*{Status}
reviewed

\section*{Safety}
safe

\section*{Shell}
posix

\section*{Platforms}
linux, gnu-linux, freebsd, openbsd, netbsd

\section*{Created At}
2026-04-25

\section*{Updated At}
2026-04-25

\end{document}
